\subsection{Schur's Lemma}

PROPOSITION 2. --- Let A be a ring, M and N two A-modules, and f a nonzero homomorphism from M to N.

\begin{enumerate}
\def\labelenumi{\alph{enumi})}
\item
  If M is simple, then f is injective.
\item
  If \(\mathbf{N}\) is simple, then \(f\) is surjective.
\item
  If M and N are simple, then f is an isomorphism.
\end{enumerate}

The kernel of f is a proper submodule of M, and the image of f is a nonzero submodule of N.

\begin{enumerate}
\def\labelenumi{\alph{enumi})}
\item
  If M is simple, then we have \(\operatorname{Ker}(f)=0\) , so f is injective.
\item
  If \(N\) is simple, then we have \(\operatorname{Im}(f) = N\) , so \(f\) is surjective.
\item
  If M and N are simple, then f is both injective and surjective.
\end{enumerate}

COROLLARY (Schur's lemma). --- The endomorphism ring of a simple module is a field.

If M is a simple A-module, then every nonzero element of the nonzero ring \(\operatorname{End}_{\mathrm{A}}(\mathrm{M})\) is invertible (Proposition 2, c)), so \(\operatorname{End}_{\mathrm{A}}(\mathrm{M})\) is a field.

THEOREM 1. --- Let K be an algebraically closed commutative field, A a K-algebra, and M a simple A-module. Suppose that the dimension of M as a vector space over K is finite or, more generally, strictly less than the cardinal of K. Then the endomorphism ring of the A-module M consists of the homotheties \(\alpha_{M}\) with \(\alpha \in K\) .

Let \(\mathrm{E}\) be the endomorphism ring of the A-module \(\mathrm{M}\) ; it is a field by the corollary of Proposition 2 and an algebra over the field \(\mathrm{K}\) . If we view \(\mathrm{M}\) as a left vector space over the field \(\mathrm{E}\) , then we have \(\dim_{\mathrm{K}} \mathrm{M} = (\dim_{\mathrm{E}} \mathrm{M})[\mathrm{E} : \mathrm{K}]\) by Proposition 25 of II, §1, No.~13, p.~222, hence \(\dim_{\mathrm{K}} \mathrm{M} \geqslant [\mathrm{E} : \mathrm{K}]\) . Since \(\dim_{\mathrm{K}} \mathrm{M} < \operatorname{Card}(\mathrm{K})\) by assumption, the equality \(\mathrm{E} = \mathrm{K} \cdot 1_{\mathrm{M}}\) is then a consequence of the following lemma.

Lemma 1. --- Let E be a field and K a subfield of the center of E not equal to E. If the field K is algebraically closed, then we have \([E:K] \geqslant \text{Card}(K)\) .

Let \(x\) be an element of \(\mathrm{E} - \mathrm{K}\) and \(\mathrm{L}\) the (commutative) subfield of \(\mathrm{E}\) generated by \(\mathrm{K} \cup \{x\}\) . Since \(\mathrm{K}\) is algebraically closed, \(x\) is transcendent over \(\mathrm{K}\) . By VII, §2, No.~3, p.~10, Theorem 2 and p.~11, the elements \((x - \alpha)^{-1}\) of \(\mathrm{L}\) , where \(\alpha\) runs through \(\mathrm{K}\) , are linearly independent over \(\mathrm{K}\) . We therefore have \([\mathrm{E} : \mathrm{K}] \geqslant [\mathrm{L} : \mathrm{K}] \geqslant \operatorname{Card}(\mathrm{K})\) .

Example. --- *Let A be a C-algebra generated by a countable family of elements; it has countable dimension over C. Let M be a simple A-module;

it is monogenous, so admits a countable basis over \(\mathbf{C}\) . Since the field \(\mathbf{C}\) is not countable (Gen.~Top., IV, §4, No.~1, p.~44), we have \([\mathrm{M}:\mathbf{C}] < \mathrm{Card}(\mathbf{C})\) . Therefore the endomorphisms of the A-module M are the homotheties \(\alpha_{\mathrm{M}}\) with \(\alpha \in \mathbf{C}\) . This applies, in particular, when A is the universal enveloping algebra of a finite-dimensional Lie algebra over \(\mathbf{C}\) (Lie, I, §2, No.~7, p.~21, Corollary 3).*

COROLLARY 1. --- Keep the assumptions of Theorem 1, and suppose, moreover, that the algebra A is commutative. Then M has dimension 1 over K.

Since the ring A is commutative, \(a_{M}\) is an endomorphism of the A-module M for every \(a \in A\) . By Theorem 1, we therefore have \(A_{M} = K \cdot 1_{M}\) , and M is a simple K-module, that is, a 1-dimensional vector space over the field K.

COROLLARY 2. --- Let K be a commutative field, A a K-algebra, and M an A-module. Suppose that for every extension L of K, the \(\mathrm{A}_{(\mathrm{L})}\) -module \(\mathrm{M}_{(\mathrm{L})}\) is simple. Then the endomorphism ring of M consists of the homotheties \(\alpha_{M}\) with \(\alpha \in K\) .

Let I be a set of cardinal strictly greater than the dimension of M over K (for example, the set of subsets of M). Let L be an algebraic closure of the field \(\mathrm{K}((\mathrm{X}_{i})_{i\in\mathrm{I}})\) (V, §4, No.~3, p.~23, Theorem 2). Choose a K-linear form \(\varphi\) on L such that \(\varphi(1)=1\) , and let \(v:M_{(\mathrm{L})}\to M\) be the K-linear mapping characterized by \(v(\alpha\otimes m)=\varphi(\alpha)m\) . Let u be an endomorphism of M. The dimension of \(M_{(\mathrm{L})}\) over L is equal to that of M over K and is strictly less than the cardinal of L. By Theorem 1, the endomorphism \(1_{L}\otimes u\) of the \(A_{(\mathrm{L})}\) -module \(M_{(\mathrm{L})}\) is of the form \(\lambda\otimes1_{M}\) with \(\lambda\in L\) . For every \(x\in M\) , we have

\[
\begin{array}{c} u (x) = v \bigl (1 \otimes u (x) \bigr) = v \bigl ((1 _ {\mathrm{L}} \otimes u) (1 \otimes x) \bigr) \\ = v \bigl ((\lambda \otimes 1 _ {\mathrm{M}}) (1 \otimes x) \bigr) = v (\lambda \otimes x) = \varphi (\lambda) x, \end{array}
\]

so that u is the homothety \(\varphi(\lambda)_{\mathrm{M}}\) .

